% paper46-semantic-futures.tex — Semantic Futures: Predictive Verification
%   via Trend Extrapolation
% Paper 46 of the Judgment Geometry series.
% Compile: pdflatex -interaction=nonstopmode paper46-semantic-futures.tex
\documentclass[11pt]{article}
\usepackage{lmodern}
% jugeo-common.tex — shared preamble for all Judgment Geometry papers
% Include via: % jugeo-common.tex — shared preamble for all Judgment Geometry papers
% Include via: % jugeo-common.tex — shared preamble for all Judgment Geometry papers
% Include via: \input{jugeo-common}

% ── Packages ────────────────────────────────────────────────────────────
\usepackage{amsmath,amssymb,amsthm,mathtools}
\usepackage{tikz-cd}
\usepackage{listings}
\usepackage{xcolor}
\usepackage{xspace}
\usepackage{booktabs}
\usepackage{multirow}
\usepackage{enumitem}
\usepackage[expansion=false]{microtype}
\usepackage{hyperref}
\usepackage{cleveref}
% \usepackage{thmtools}  % disabled: not available in this TeX installation
\usepackage{float}
\usepackage{caption}
\usepackage{subcaption}
\usepackage{tabularx}
\usepackage{stmaryrd}       % \llbracket, \rrbracket
\usepackage{bbm}            % \mathbbm{1}

% ── Page geometry ───────────────────────────────────────────────────────
\usepackage[margin=1in]{geometry}

% ── Theorem environments ────────────────────────────────────────────────
\theoremstyle{plain}
\newtheorem{theorem}{Theorem}[section]
\newtheorem{lemma}[theorem]{Lemma}
\newtheorem{proposition}[theorem]{Proposition}
\newtheorem{corollary}[theorem]{Corollary}

\theoremstyle{definition}
\newtheorem{definition}[theorem]{Definition}
\newtheorem{example}[theorem]{Example}
\newtheorem{construction}[theorem]{Construction}

\theoremstyle{remark}
\newtheorem{remark}[theorem]{Remark}
\newtheorem{notation}[theorem]{Notation}

% ── Python listings style ──────────────────────────────────────────────
\definecolor{jugeoBlue}{HTML}{2563EB}
\definecolor{jugeoGreen}{HTML}{059669}
\definecolor{jugeoRed}{HTML}{DC2626}
\definecolor{jugeoPurple}{HTML}{7C3AED}
\definecolor{jugeoGray}{HTML}{6B7280}
\definecolor{jugeoBg}{HTML}{F8FAFC}

\lstdefinestyle{jugeo-python}{
  language=Python,
  basicstyle=\ttfamily\small,
  keywordstyle=\color{jugeoBlue}\bfseries,
  stringstyle=\color{jugeoGreen},
  commentstyle=\color{jugeoGray}\itshape,
  numberstyle=\tiny\color{jugeoGray},
  backgroundcolor=\color{jugeoBg},
  numbers=left,
  numbersep=6pt,
  frame=single,
  framerule=0.4pt,
  rulecolor=\color{jugeoGray!40},
  breaklines=true,
  breakatwhitespace=true,
  showstringspaces=false,
  tabsize=4,
  xleftmargin=1.5em,
  framexleftmargin=1.5em,
  morekeywords={dataclass,frozen,field,Enum,IntEnum,auto,Any,
                Mapping,Sequence,Optional,match,case},
  literate={→}{{$\to$}}1 {←}{{$\leftarrow$}}1
           {≤}{{$\leq$}}1 {≥}{{$\geq$}}1 {≠}{{$\neq$}}1
           {∈}{{$\in$}}1 {∀}{{$\forall$}}1 {∃}{{$\exists$}}1
           {⊕}{{$\oplus$}}1 {⊖}{{$\ominus$}}1,
  escapeinside={(*@}{@*)},
}
\lstset{style=jugeo-python}

\lstdefinestyle{jugeo-output}{
  basicstyle=\ttfamily\small,
  backgroundcolor=\color{jugeoBg},
  frame=single,
  framerule=0.4pt,
  rulecolor=\color{jugeoGray!40},
  breaklines=true,
  showstringspaces=false,
  xleftmargin=1.5em,
  framexleftmargin=0.5em,
  numbers=none,
}

% ── JuGeo-specific macros ──────────────────────────────────────────────

% Judgment 8-tuple
\newcommand{\J}{\mathcal{J}}
\newcommand{\Jud}[8]{(#1,\,#2,\,#3,\,#4,\,#5,\,#6,\,#7,\,#8)}
\newcommand{\JudShort}[1]{\J_{#1}}

% Components
\newcommand{\coord}{\mathsf{c}}            % coordinate
\newcommand{\prop}{\varphi}                 % proposition
\newcommand{\carrier}{\mathcal{A}}          % carrier type
\newcommand{\evid}{\mathcal{E}}             % evidence bundle
\newcommand{\oblig}{\mathcal{O}}            % obligations
\newcommand{\obstr}{\mathcal{B}}            % obstructions
\newcommand{\trust}{\mathcal{T}}            % trust annotation
\newcommand{\prov}{\Pi}                     % provenance

% Trust algebra
\newcommand{\Talg}{\mathfrak{T}}            % trust ordered algebra
\newcommand{\Eadm}{\mathcal{E}_{\mathrm{adm}}}
\newcommand{\tleq}{\preceq}                 % trust ordering
\newcommand{\tjoin}{\oplus}                 % conservative join
\newcommand{\tatten}{\ominus}               % attenuation
\newcommand{\tpromote}{\uparrow_\pi}        % promotion
\newcommand{\tdemote}{\downarrow_\chi}       % demotion

% Trust tiers
\newcommand{\Tcontra}{\ensuremath{\mathsf{CONTRADICTED}}}
\newcommand{\Tunver}{\ensuremath{\mathsf{UNVERIFIED}}}
\newcommand{\Tcopilot}{\ensuremath{\mathsf{COPILOT}}}
\newcommand{\Toracle}{\ensuremath{\mathsf{ORACLE}}}
\newcommand{\Truntime}{\ensuremath{\mathsf{RUNTIME}}}
\newcommand{\Tsolver}{\ensuremath{\mathsf{SOLVER}}}
\newcommand{\Tproof}{\ensuremath{\mathsf{PROOF}}}

% Site and geometry
\newcommand{\Site}{\mathcal{S}}             % semantic site
\newcommand{\Cov}{\mathrm{Cov}}            % covering family
\newcommand{\Ob}{\mathrm{Ob}}              % objects
\newcommand{\Mor}{\mathrm{Mor}}            % morphisms
\newcommand{\res}{\mathrm{res}}            % restriction
\newcommand{\Cech}{\check{C}}              % Čech complex
\newcommand{\Hcoh}[1]{H^{#1}}             % cohomology group

% Descent
\newcommand{\descent}{\mathrm{desc}}
\newcommand{\glue}{\mathrm{glue}}
\newcommand{\overlap}[2]{#1 \cap #2}

% Semantic moves
\newcommand{\VERIFY}{\textsc{Verify}}
\newcommand{\CONSTRUCT}{\textsc{Construct}}
\newcommand{\NORMALIZE}{\textsc{Normalize}}
\newcommand{\GLUE}{\textsc{Glue}}
\newcommand{\RESTRICT}{\textsc{Restrict}}
\newcommand{\TRANSPORT}{\textsc{Transport}}
\newcommand{\REPAIR}{\textsc{Repair}}
\newcommand{\PROMOTE}{\textsc{Promote}}
\newcommand{\CHALLENGE}{\textsc{Challenge}}
\newcommand{\ESCALATE}{\textsc{Escalate}}

% Fragments
\newcommand{\QFLIA}{\texttt{QF\_LIA}}
\newcommand{\QFLRA}{\texttt{QF\_LRA}}
\newcommand{\QFBV}{\texttt{QF\_BV}}
\newcommand{\QFUF}{\texttt{QF\_UF}}

% Misc
\newcommand{\Python}{\textsc{Python}}
\newcommand{\JuGeo}{\textsc{JuGeo}}
\newcommand{\LEAN}{\textsc{Lean}\,4}
\newcommand{\Fstar}{F$^{\star}$}
\newcommand{\Coq}{\textsc{Coq}}
\newcommand{\Dafny}{\textsc{Dafny}}

% Common shortcuts
\newcommand{\ie}{i.e.\@\xspace}
\newcommand{\eg}{e.g.\@\xspace}
\newcommand{\cf}{cf.\@\xspace}
\newcommand{\wrt}{w.r.t.\@\xspace}

% Evaluation shortcuts
\newcommand{\numcases}{300}
\newcommand{\numspec}{100}
\newcommand{\numeq}{100}
\newcommand{\numbug}{100}
\newcommand{\medtime}{6.5\,ms}
\newcommand{\totaltime}{1.949\,s}


% ── Packages ────────────────────────────────────────────────────────────
\usepackage{amsmath,amssymb,amsthm,mathtools}
\usepackage{tikz-cd}
\usepackage{listings}
\usepackage{xcolor}
\usepackage{xspace}
\usepackage{booktabs}
\usepackage{multirow}
\usepackage{enumitem}
\usepackage[expansion=false]{microtype}
\usepackage{hyperref}
\usepackage{cleveref}
% \usepackage{thmtools}  % disabled: not available in this TeX installation
\usepackage{float}
\usepackage{caption}
\usepackage{subcaption}
\usepackage{tabularx}
\usepackage{stmaryrd}       % \llbracket, \rrbracket
\usepackage{bbm}            % \mathbbm{1}

% ── Page geometry ───────────────────────────────────────────────────────
\usepackage[margin=1in]{geometry}

% ── Theorem environments ────────────────────────────────────────────────
\theoremstyle{plain}
\newtheorem{theorem}{Theorem}[section]
\newtheorem{lemma}[theorem]{Lemma}
\newtheorem{proposition}[theorem]{Proposition}
\newtheorem{corollary}[theorem]{Corollary}

\theoremstyle{definition}
\newtheorem{definition}[theorem]{Definition}
\newtheorem{example}[theorem]{Example}
\newtheorem{construction}[theorem]{Construction}

\theoremstyle{remark}
\newtheorem{remark}[theorem]{Remark}
\newtheorem{notation}[theorem]{Notation}

% ── Python listings style ──────────────────────────────────────────────
\definecolor{jugeoBlue}{HTML}{2563EB}
\definecolor{jugeoGreen}{HTML}{059669}
\definecolor{jugeoRed}{HTML}{DC2626}
\definecolor{jugeoPurple}{HTML}{7C3AED}
\definecolor{jugeoGray}{HTML}{6B7280}
\definecolor{jugeoBg}{HTML}{F8FAFC}

\lstdefinestyle{jugeo-python}{
  language=Python,
  basicstyle=\ttfamily\small,
  keywordstyle=\color{jugeoBlue}\bfseries,
  stringstyle=\color{jugeoGreen},
  commentstyle=\color{jugeoGray}\itshape,
  numberstyle=\tiny\color{jugeoGray},
  backgroundcolor=\color{jugeoBg},
  numbers=left,
  numbersep=6pt,
  frame=single,
  framerule=0.4pt,
  rulecolor=\color{jugeoGray!40},
  breaklines=true,
  breakatwhitespace=true,
  showstringspaces=false,
  tabsize=4,
  xleftmargin=1.5em,
  framexleftmargin=1.5em,
  morekeywords={dataclass,frozen,field,Enum,IntEnum,auto,Any,
                Mapping,Sequence,Optional,match,case},
  literate={→}{{$\to$}}1 {←}{{$\leftarrow$}}1
           {≤}{{$\leq$}}1 {≥}{{$\geq$}}1 {≠}{{$\neq$}}1
           {∈}{{$\in$}}1 {∀}{{$\forall$}}1 {∃}{{$\exists$}}1
           {⊕}{{$\oplus$}}1 {⊖}{{$\ominus$}}1,
  escapeinside={(*@}{@*)},
}
\lstset{style=jugeo-python}

\lstdefinestyle{jugeo-output}{
  basicstyle=\ttfamily\small,
  backgroundcolor=\color{jugeoBg},
  frame=single,
  framerule=0.4pt,
  rulecolor=\color{jugeoGray!40},
  breaklines=true,
  showstringspaces=false,
  xleftmargin=1.5em,
  framexleftmargin=0.5em,
  numbers=none,
}

% ── JuGeo-specific macros ──────────────────────────────────────────────

% Judgment 8-tuple
\newcommand{\J}{\mathcal{J}}
\newcommand{\Jud}[8]{(#1,\,#2,\,#3,\,#4,\,#5,\,#6,\,#7,\,#8)}
\newcommand{\JudShort}[1]{\J_{#1}}

% Components
\newcommand{\coord}{\mathsf{c}}            % coordinate
\newcommand{\prop}{\varphi}                 % proposition
\newcommand{\carrier}{\mathcal{A}}          % carrier type
\newcommand{\evid}{\mathcal{E}}             % evidence bundle
\newcommand{\oblig}{\mathcal{O}}            % obligations
\newcommand{\obstr}{\mathcal{B}}            % obstructions
\newcommand{\trust}{\mathcal{T}}            % trust annotation
\newcommand{\prov}{\Pi}                     % provenance

% Trust algebra
\newcommand{\Talg}{\mathfrak{T}}            % trust ordered algebra
\newcommand{\Eadm}{\mathcal{E}_{\mathrm{adm}}}
\newcommand{\tleq}{\preceq}                 % trust ordering
\newcommand{\tjoin}{\oplus}                 % conservative join
\newcommand{\tatten}{\ominus}               % attenuation
\newcommand{\tpromote}{\uparrow_\pi}        % promotion
\newcommand{\tdemote}{\downarrow_\chi}       % demotion

% Trust tiers
\newcommand{\Tcontra}{\ensuremath{\mathsf{CONTRADICTED}}}
\newcommand{\Tunver}{\ensuremath{\mathsf{UNVERIFIED}}}
\newcommand{\Tcopilot}{\ensuremath{\mathsf{COPILOT}}}
\newcommand{\Toracle}{\ensuremath{\mathsf{ORACLE}}}
\newcommand{\Truntime}{\ensuremath{\mathsf{RUNTIME}}}
\newcommand{\Tsolver}{\ensuremath{\mathsf{SOLVER}}}
\newcommand{\Tproof}{\ensuremath{\mathsf{PROOF}}}

% Site and geometry
\newcommand{\Site}{\mathcal{S}}             % semantic site
\newcommand{\Cov}{\mathrm{Cov}}            % covering family
\newcommand{\Ob}{\mathrm{Ob}}              % objects
\newcommand{\Mor}{\mathrm{Mor}}            % morphisms
\newcommand{\res}{\mathrm{res}}            % restriction
\newcommand{\Cech}{\check{C}}              % Čech complex
\newcommand{\Hcoh}[1]{H^{#1}}             % cohomology group

% Descent
\newcommand{\descent}{\mathrm{desc}}
\newcommand{\glue}{\mathrm{glue}}
\newcommand{\overlap}[2]{#1 \cap #2}

% Semantic moves
\newcommand{\VERIFY}{\textsc{Verify}}
\newcommand{\CONSTRUCT}{\textsc{Construct}}
\newcommand{\NORMALIZE}{\textsc{Normalize}}
\newcommand{\GLUE}{\textsc{Glue}}
\newcommand{\RESTRICT}{\textsc{Restrict}}
\newcommand{\TRANSPORT}{\textsc{Transport}}
\newcommand{\REPAIR}{\textsc{Repair}}
\newcommand{\PROMOTE}{\textsc{Promote}}
\newcommand{\CHALLENGE}{\textsc{Challenge}}
\newcommand{\ESCALATE}{\textsc{Escalate}}

% Fragments
\newcommand{\QFLIA}{\texttt{QF\_LIA}}
\newcommand{\QFLRA}{\texttt{QF\_LRA}}
\newcommand{\QFBV}{\texttt{QF\_BV}}
\newcommand{\QFUF}{\texttt{QF\_UF}}

% Misc
\newcommand{\Python}{\textsc{Python}}
\newcommand{\JuGeo}{\textsc{JuGeo}}
\newcommand{\LEAN}{\textsc{Lean}\,4}
\newcommand{\Fstar}{F$^{\star}$}
\newcommand{\Coq}{\textsc{Coq}}
\newcommand{\Dafny}{\textsc{Dafny}}

% Common shortcuts
\newcommand{\ie}{i.e.\@\xspace}
\newcommand{\eg}{e.g.\@\xspace}
\newcommand{\cf}{cf.\@\xspace}
\newcommand{\wrt}{w.r.t.\@\xspace}

% Evaluation shortcuts
\newcommand{\numcases}{300}
\newcommand{\numspec}{100}
\newcommand{\numeq}{100}
\newcommand{\numbug}{100}
\newcommand{\medtime}{6.5\,ms}
\newcommand{\totaltime}{1.949\,s}


% ── Packages ────────────────────────────────────────────────────────────
\usepackage{amsmath,amssymb,amsthm,mathtools}
\usepackage{tikz-cd}
\usepackage{listings}
\usepackage{xcolor}
\usepackage{xspace}
\usepackage{booktabs}
\usepackage{multirow}
\usepackage{enumitem}
\usepackage[expansion=false]{microtype}
\usepackage{hyperref}
\usepackage{cleveref}
% \usepackage{thmtools}  % disabled: not available in this TeX installation
\usepackage{float}
\usepackage{caption}
\usepackage{subcaption}
\usepackage{tabularx}
\usepackage{stmaryrd}       % \llbracket, \rrbracket
\usepackage{bbm}            % \mathbbm{1}

% ── Page geometry ───────────────────────────────────────────────────────
\usepackage[margin=1in]{geometry}

% ── Theorem environments ────────────────────────────────────────────────
\theoremstyle{plain}
\newtheorem{theorem}{Theorem}[section]
\newtheorem{lemma}[theorem]{Lemma}
\newtheorem{proposition}[theorem]{Proposition}
\newtheorem{corollary}[theorem]{Corollary}

\theoremstyle{definition}
\newtheorem{definition}[theorem]{Definition}
\newtheorem{example}[theorem]{Example}
\newtheorem{construction}[theorem]{Construction}

\theoremstyle{remark}
\newtheorem{remark}[theorem]{Remark}
\newtheorem{notation}[theorem]{Notation}

% ── Python listings style ──────────────────────────────────────────────
\definecolor{jugeoBlue}{HTML}{2563EB}
\definecolor{jugeoGreen}{HTML}{059669}
\definecolor{jugeoRed}{HTML}{DC2626}
\definecolor{jugeoPurple}{HTML}{7C3AED}
\definecolor{jugeoGray}{HTML}{6B7280}
\definecolor{jugeoBg}{HTML}{F8FAFC}

\lstdefinestyle{jugeo-python}{
  language=Python,
  basicstyle=\ttfamily\small,
  keywordstyle=\color{jugeoBlue}\bfseries,
  stringstyle=\color{jugeoGreen},
  commentstyle=\color{jugeoGray}\itshape,
  numberstyle=\tiny\color{jugeoGray},
  backgroundcolor=\color{jugeoBg},
  numbers=left,
  numbersep=6pt,
  frame=single,
  framerule=0.4pt,
  rulecolor=\color{jugeoGray!40},
  breaklines=true,
  breakatwhitespace=true,
  showstringspaces=false,
  tabsize=4,
  xleftmargin=1.5em,
  framexleftmargin=1.5em,
  morekeywords={dataclass,frozen,field,Enum,IntEnum,auto,Any,
                Mapping,Sequence,Optional,match,case},
  literate={→}{{$\to$}}1 {←}{{$\leftarrow$}}1
           {≤}{{$\leq$}}1 {≥}{{$\geq$}}1 {≠}{{$\neq$}}1
           {∈}{{$\in$}}1 {∀}{{$\forall$}}1 {∃}{{$\exists$}}1
           {⊕}{{$\oplus$}}1 {⊖}{{$\ominus$}}1,
  escapeinside={(*@}{@*)},
}
\lstset{style=jugeo-python}

\lstdefinestyle{jugeo-output}{
  basicstyle=\ttfamily\small,
  backgroundcolor=\color{jugeoBg},
  frame=single,
  framerule=0.4pt,
  rulecolor=\color{jugeoGray!40},
  breaklines=true,
  showstringspaces=false,
  xleftmargin=1.5em,
  framexleftmargin=0.5em,
  numbers=none,
}

% ── JuGeo-specific macros ──────────────────────────────────────────────

% Judgment 8-tuple
\newcommand{\J}{\mathcal{J}}
\newcommand{\Jud}[8]{(#1,\,#2,\,#3,\,#4,\,#5,\,#6,\,#7,\,#8)}
\newcommand{\JudShort}[1]{\J_{#1}}

% Components
\newcommand{\coord}{\mathsf{c}}            % coordinate
\newcommand{\prop}{\varphi}                 % proposition
\newcommand{\carrier}{\mathcal{A}}          % carrier type
\newcommand{\evid}{\mathcal{E}}             % evidence bundle
\newcommand{\oblig}{\mathcal{O}}            % obligations
\newcommand{\obstr}{\mathcal{B}}            % obstructions
\newcommand{\trust}{\mathcal{T}}            % trust annotation
\newcommand{\prov}{\Pi}                     % provenance

% Trust algebra
\newcommand{\Talg}{\mathfrak{T}}            % trust ordered algebra
\newcommand{\Eadm}{\mathcal{E}_{\mathrm{adm}}}
\newcommand{\tleq}{\preceq}                 % trust ordering
\newcommand{\tjoin}{\oplus}                 % conservative join
\newcommand{\tatten}{\ominus}               % attenuation
\newcommand{\tpromote}{\uparrow_\pi}        % promotion
\newcommand{\tdemote}{\downarrow_\chi}       % demotion

% Trust tiers
\newcommand{\Tcontra}{\ensuremath{\mathsf{CONTRADICTED}}}
\newcommand{\Tunver}{\ensuremath{\mathsf{UNVERIFIED}}}
\newcommand{\Tcopilot}{\ensuremath{\mathsf{COPILOT}}}
\newcommand{\Toracle}{\ensuremath{\mathsf{ORACLE}}}
\newcommand{\Truntime}{\ensuremath{\mathsf{RUNTIME}}}
\newcommand{\Tsolver}{\ensuremath{\mathsf{SOLVER}}}
\newcommand{\Tproof}{\ensuremath{\mathsf{PROOF}}}

% Site and geometry
\newcommand{\Site}{\mathcal{S}}             % semantic site
\newcommand{\Cov}{\mathrm{Cov}}            % covering family
\newcommand{\Ob}{\mathrm{Ob}}              % objects
\newcommand{\Mor}{\mathrm{Mor}}            % morphisms
\newcommand{\res}{\mathrm{res}}            % restriction
\newcommand{\Cech}{\check{C}}              % Čech complex
\newcommand{\Hcoh}[1]{H^{#1}}             % cohomology group

% Descent
\newcommand{\descent}{\mathrm{desc}}
\newcommand{\glue}{\mathrm{glue}}
\newcommand{\overlap}[2]{#1 \cap #2}

% Semantic moves
\newcommand{\VERIFY}{\textsc{Verify}}
\newcommand{\CONSTRUCT}{\textsc{Construct}}
\newcommand{\NORMALIZE}{\textsc{Normalize}}
\newcommand{\GLUE}{\textsc{Glue}}
\newcommand{\RESTRICT}{\textsc{Restrict}}
\newcommand{\TRANSPORT}{\textsc{Transport}}
\newcommand{\REPAIR}{\textsc{Repair}}
\newcommand{\PROMOTE}{\textsc{Promote}}
\newcommand{\CHALLENGE}{\textsc{Challenge}}
\newcommand{\ESCALATE}{\textsc{Escalate}}

% Fragments
\newcommand{\QFLIA}{\texttt{QF\_LIA}}
\newcommand{\QFLRA}{\texttt{QF\_LRA}}
\newcommand{\QFBV}{\texttt{QF\_BV}}
\newcommand{\QFUF}{\texttt{QF\_UF}}

% Misc
\newcommand{\Python}{\textsc{Python}}
\newcommand{\JuGeo}{\textsc{JuGeo}}
\newcommand{\LEAN}{\textsc{Lean}\,4}
\newcommand{\Fstar}{F$^{\star}$}
\newcommand{\Coq}{\textsc{Coq}}
\newcommand{\Dafny}{\textsc{Dafny}}

% Common shortcuts
\newcommand{\ie}{i.e.\@\xspace}
\newcommand{\eg}{e.g.\@\xspace}
\newcommand{\cf}{cf.\@\xspace}
\newcommand{\wrt}{w.r.t.\@\xspace}

% Evaluation shortcuts
\newcommand{\numcases}{300}
\newcommand{\numspec}{100}
\newcommand{\numeq}{100}
\newcommand{\numbug}{100}
\newcommand{\medtime}{6.5\,ms}
\newcommand{\totaltime}{1.949\,s}

% experiment-data.tex — AUTO-GENERATED from real jugeo CLI runs
% DO NOT EDIT — regenerate with: python3 experiments/run_universal_metrics.py
% Generated from 30 programs across 6 domains

\newcommand{\expTotalPrograms}{30}
\newcommand{\expVerified}{30}
\newcommand{\expAccuracy}{100.0\%}
\newcommand{\expCoordsMin}{2}
\newcommand{\expCoordsMax}{10}
\newcommand{\expCoordsMean}{5.2}
\newcommand{\expCoordsMedian}{5.5}
\newcommand{\expPropsMin}{4}
\newcommand{\expPropsMax}{20}
\newcommand{\expPropsMean}{10.4}
\newcommand{\expPropsSum}{311}
\newcommand{\expPropsOkSum}{310}
\newcommand{\expTimeMin}{3.506\,s}
\newcommand{\expTimeMax}{6.714\,s}
\newcommand{\expTimeMean}{4.64\,s}
\newcommand{\expTimeMedian}{4.508\,s}
\newcommand{\expTimeTotal}{139.204\,s}
\newcommand{\expObstructionTotal}{0}
\newcommand{\expHOneZero}{30}
\newcommand{\expDomainCount}{6}
\newcommand{\expTrustCOPILOTSUGGESTED}{30}
\newcommand{\expDomainDsCount}{5}
\newcommand{\expDomainDsVerified}{5}
\newcommand{\expDomainDsAvgCoords}{7.6}
\newcommand{\expDomainDsAvgProps}{15.2}
\newcommand{\expDomainMathCount}{5}
\newcommand{\expDomainMathVerified}{5}
\newcommand{\expDomainMathAvgCoords}{4.8}
\newcommand{\expDomainMathAvgProps}{9.4}
\newcommand{\expDomainSmCount}{5}
\newcommand{\expDomainSmVerified}{5}
\newcommand{\expDomainSmAvgCoords}{7.6}
\newcommand{\expDomainSmAvgProps}{15.2}
\newcommand{\expDomainSortCount}{5}
\newcommand{\expDomainSortVerified}{5}
\newcommand{\expDomainSortAvgCoords}{2.4}
\newcommand{\expDomainSortAvgProps}{4.8}
\newcommand{\expDomainStrCount}{5}
\newcommand{\expDomainStrVerified}{5}
\newcommand{\expDomainStrAvgCoords}{3.2}
\newcommand{\expDomainStrAvgProps}{6.4}
\newcommand{\expDomainWebCount}{5}
\newcommand{\expDomainWebVerified}{5}
\newcommand{\expDomainWebAvgCoords}{5.6}
\newcommand{\expDomainWebAvgProps}{11.2}

% subsystem-data.tex — AUTO-GENERATED from real jugeo subsystem experiments
% DO NOT EDIT — regenerate with: python3 experiments/run_subsystem_experiments.py

\newcommand{\subCliTotal}{10}
\newcommand{\subCliVerified}{9}
\newcommand{\subCliAccuracy}{90.0\%}
\newcommand{\subCliMeanTime}{4.132\,s}
\newcommand{\subCliMinTime}{3.472\,s}
\newcommand{\subCliMaxTime}{5.372\,s}
\newcommand{\subCliTotalCoords}{54}
\newcommand{\subCliTotalProps}{107}
\newcommand{\subCliTotalPropsOk}{106}
\newcommand{\subSiteCalls}{90}
\newcommand{\subSiteSuccessful}{69}
\newcommand{\subSiteSuccessRate}{76.7\%}
\newcommand{\subSiteMeanTime}{0.0001\,s}
\newcommand{\subSiteTotalTime}{0.005\,s}
\newcommand{\subSpecificationsatisfactionSmallTime}{0.0003\,s}
\newcommand{\subSpecificationsatisfactionMediumTime}{0.0001\,s}
\newcommand{\subSpecificationsatisfactionLargeTime}{0.0001\,s}
\newcommand{\subInhabitantfleetSmallTime}{0.0\,s}
\newcommand{\subInhabitantfleetMediumTime}{0.0\,s}
\newcommand{\subInhabitantfleetLargeTime}{0.0\,s}
\newcommand{\subTheoremecologySmallTime}{0.0\,s}
\newcommand{\subTheoremecologyMediumTime}{0.0\,s}
\newcommand{\subTheoremecologyLargeTime}{0.0\,s}
\newcommand{\subStatespaceexplorationSmallTime}{0.0\,s}
\newcommand{\subStatespaceexplorationMediumTime}{0.0\,s}
\newcommand{\subStatespaceexplorationLargeTime}{0.0\,s}
\newcommand{\subRepairsemanticsSmallTime}{0.0\,s}
\newcommand{\subRepairsemanticsMediumTime}{0.0\,s}
\newcommand{\subRepairsemanticsLargeTime}{0.0\,s}
\newcommand{\subReplaygluingSmallTime}{0.0\,s}
\newcommand{\subReplaygluingMediumTime}{0.0\,s}
\newcommand{\subReplaygluingLargeTime}{0.0\,s}
\newcommand{\subSemanticfuturesSmallTime}{0.0\,s}
\newcommand{\subSemanticfuturesMediumTime}{0.0\,s}
\newcommand{\subSemanticfuturesLargeTime}{0.0\,s}
\newcommand{\subHypercovertreatySmallTime}{0.0\,s}
\newcommand{\subHypercovertreatyMediumTime}{0.0\,s}
\newcommand{\subHypercovertreatyLargeTime}{0.0\,s}
\newcommand{\subEncodeforsolverSmallTime}{0.0026\,s}
\newcommand{\subEncodeforsolverMediumTime}{0.0002\,s}
\newcommand{\subEncodeforsolverLargeTime}{0.0001\,s}
\newcommand{\subInterfaceroutingSmallTime}{0.0\,s}
\newcommand{\subInterfaceroutingMediumTime}{0.0\,s}
\newcommand{\subInterfaceroutingLargeTime}{0.0\,s}
\newcommand{\subOrchestrateverificationSmallTime}{0.0002\,s}
\newcommand{\subOrchestrateverificationMediumTime}{0.0\,s}
\newcommand{\subOrchestrateverificationLargeTime}{0.0\,s}
\newcommand{\subRunfulldescentSmallTime}{0.0\,s}
\newcommand{\subRunfulldescentMediumTime}{0.0\,s}
\newcommand{\subRunfulldescentLargeTime}{0.0\,s}
\newcommand{\subKernellifecycleSmallTime}{0.0\,s}
\newcommand{\subKernellifecycleMediumTime}{0.0\,s}
\newcommand{\subKernellifecycleLargeTime}{0.0\,s}
\newcommand{\subDiscoverypipelineSmallTime}{0.0\,s}
\newcommand{\subDiscoverypipelineMediumTime}{0.0\,s}
\newcommand{\subDiscoverypipelineLargeTime}{0.0\,s}
\newcommand{\subAnalogytransportSmallTime}{0.0\,s}
\newcommand{\subAnalogytransportMediumTime}{0.0\,s}
\newcommand{\subAnalogytransportLargeTime}{0.0\,s}
\newcommand{\subFormalcoresiteSmallTime}{0.0001\,s}
\newcommand{\subFormalcoresiteMediumTime}{0.0\,s}
\newcommand{\subFormalcoresiteLargeTime}{0.0\,s}
\newcommand{\subProblematlasSmallTime}{0.0\,s}
\newcommand{\subProblematlasMediumTime}{0.0\,s}
\newcommand{\subProblematlasLargeTime}{0.0\,s}
\newcommand{\subPublicalignmentSmallTime}{0.0\,s}
\newcommand{\subPublicalignmentMediumTime}{0.0\,s}
\newcommand{\subPublicalignmentLargeTime}{0.0\,s}
\newcommand{\subRegimebootstrappingSmallTime}{0.0\,s}
\newcommand{\subRegimebootstrappingMediumTime}{0.0\,s}
\newcommand{\subRegimebootstrappingLargeTime}{0.0\,s}
\newcommand{\subRelationalrefinementSmallTime}{0.0\,s}
\newcommand{\subRelationalrefinementMediumTime}{0.0\,s}
\newcommand{\subRelationalrefinementLargeTime}{0.0\,s}
\newcommand{\subSemanticclosureSmallTime}{0.0\,s}
\newcommand{\subSemanticclosureMediumTime}{0.0\,s}
\newcommand{\subSemanticclosureLargeTime}{0.0\,s}
\newcommand{\subTheoremeconomicsSmallTime}{0.0001\,s}
\newcommand{\subTheoremeconomicsMediumTime}{0.0\,s}
\newcommand{\subTheoremeconomicsLargeTime}{0.0\,s}
\newcommand{\subBenchmarksuiteSmallTime}{0.0\,s}
\newcommand{\subBenchmarksuiteMediumTime}{0.0\,s}
\newcommand{\subBenchmarksuiteLargeTime}{0.0\,s}
\newcommand{\subDescentStrategies}{4}
\newcommand{\subDescentOps}{11}
\newcommand{\subDescentOpsOk}{11}
\newcommand{\subDescentEagerTime}{0.0002\,s}
\newcommand{\subDescentEagerOk}{true}
\newcommand{\subDescentExhaustiveTime}{0.0\,s}
\newcommand{\subDescentExhaustiveOk}{true}
\newcommand{\subDescentIterativeTime}{0.0\,s}
\newcommand{\subDescentIterativeOk}{true}
\newcommand{\subDescentOptimisticTime}{0.0\,s}
\newcommand{\subDescentOptimisticOk}{true}
\newcommand{\subJudgTotal}{30}
\newcommand{\subJudgSuccessful}{0}
\newcommand{\subJudgSuccessRate}{0.0\%}
\newcommand{\subJudgMeanTimeUs}{7.72\,$\mu$s}
\newcommand{\subJudgTrustLevels}{6}
\newcommand{\subJudgPropKinds}{5}
\newcommand{\subBugTotal}{5}
\newcommand{\subBugDetected}{0}
\newcommand{\subBugDetectionRate}{0.0\%}
\newcommand{\subBugFalsePositiveRate}{0.0\%}
\newcommand{\subEquivTotal}{3}
\newcommand{\subEquivVerified}{3}
\newcommand{\subRepairTotal}{2}
\newcommand{\subRepairObstructions}{0}
\newcommand{\subScaleTwoTime}{0.0001\,s}
\newcommand{\subScaleFourTime}{0.0\,s}
\newcommand{\subScaleEightTime}{0.0\,s}
\newcommand{\subScaleTwelveTime}{0.0\,s}
\newcommand{\subScaleSixteenTime}{0.0\,s}
\newcommand{\subScaleTwentyTime}{0.0\,s}
\newcommand{\subTotalExperimentTime}{95.6\,s}

% data-paper46.tex -- AUTO-GENERATED by exp46_semantic_futures.py
% DO NOT EDIT -- regenerate with: python3 experiments/exp46_semantic_futures.py

\newcommand{\ppFortySixTotalPrograms}{10}
\newcommand{\ppFortySixTotalCycles}{30}

% --- Budget saved (replaces --- in table) ---
\newcommand{\ppFortySixBudgetSaved}{0.0\%}

% --- Cycle metrics ---
\newcommand{\ppFortySixMeanCycleDuration}{99\,\mu s}
\newcommand{\ppFortySixMeanTrustScore}{0.50}
\newcommand{\ppFortySixSuccessRate}{100.0\%}
\newcommand{\ppFortySixObstructionRate}{0.0\%}

% --- Prediction metrics ---
\newcommand{\ppFortySixPredictionAccuracy}{100.0\%}
\newcommand{\ppFortySixMeanStepsToDecide}{5.0}
\newcommand{\ppFortySixEarlyTermRate}{0.0\%}


% ── Additional packages ────────────────────────────────────────────────
\usepackage{array}
\usepackage{pgfplots}
\pgfplotsset{compat=1.18}

% ── Paper-specific macros ──────────────────────────────────────────────
\newcommand{\SF}{\mathsf{SF}}               % SemanticFuture
\newcommand{\SFut}[1]{\SF_{#1}}             % indexed future
\newcommand{\FPred}{\mathsf{FPred}}         % FuturePredictor
\newcommand{\FState}{\mathsf{FState}}       % FutureState
\newcommand{\Trend}{\mathsf{Trend}}         % trend direction
\newcommand{\Horizon}{\mathsf{h}}           % prediction horizon
\newcommand{\Budget}{\mathcal{B}}           % verification budget
\newcommand{\BudgetRem}[1]{\Budget_{\mathrm{rem}}(#1)}
\newcommand{\DescentSeq}{(d_n)_{n\ge 0}}    % descent sequence
\newcommand{\DescentN}{d_n}                 % descent value at n
\newcommand{\DescentRate}{r}                % convergence rate
\newcommand{\PredSucc}{\hat{s}}             % predicted success
\newcommand{\PredFail}{\hat{f}}             % predicted failure
\newcommand{\EarlyStop}{\mathsf{Stop}}      % early-stop predicate
\newcommand{\RegimeK}{\mathsf{RK}}          % RegimeKind
\newcommand{\RKcov}{\textsc{CoverRefinement}}
\newcommand{\RKtheo}{\textsc{TheoryExtension}}
\newcommand{\RKpack}{\textsc{PackChange}}
\newcommand{\RKinv}{\textsc{InvariantStrengthening}}
\newcommand{\BootReg}{\mathsf{Boot}}        % regime bootstrapping
\newcommand{\ObsRank}{\omega}               % obstruction rank
\newcommand{\HOne}{H^1}                     % first cohomology
\newcommand{\HOneProg}[1]{\HOne(#1)}        % H^1 of program
\newcommand{\Thresh}{\tau}                  % threshold parameter
\newcommand{\Predict}{\mathsf{predict}}     % prediction function
\newcommand{\WindowSize}{w}                 % sliding window size
\newcommand{\Tags}{\mathsf{Tags}}           % tag set
\newcommand{\StepsToDecide}{T^{\star}}      % decision time
\newcommand{\maxstep}{N}                    % max steps
\newcommand{\msec}{\,\mathrm{ms}}
\newcommand{\LinReg}{\mathrm{LinReg}}       % linear regression
% ─────────────────────────────────────────────────────────────────────

\title{\textbf{Semantic Futures: Predictive Verification\\
       via Trend Extrapolation}}

\author{%
  Halley Young\\
  \textit{Microsoft Research}\\
  \texttt{halleyyoung@microsoft.com}
}

\date{\today}

\begin{document}
\maketitle

\begin{abstract}
Formal verification is expensive.  When a descent-based checker
applies semantic moves to a program judgment, it may spend thousands of
milliseconds on a proof search that is, in retrospect, doomed to fail:
the local sections will never cohere into a global section because a
first-cohomology obstruction is growing, not shrinking.  We introduce
\emph{Semantic Futures}, a predictive layer implemented in
\JuGeo's \texttt{ideation/semantic\_futures/} sub-package that
extrapolates current verification trends to forecast outcomes before
the budget is exhausted.

The \texttt{SemanticFuture} data model captures the projected state of a
verification at a future step horizon, together with a trend direction
and a confidence score.  The \texttt{FuturePredictor} class maintains a
sliding window of descent values and obstruction ranks, fits a linear
trend, and classifies the trajectory as \emph{converging} (predict
success), \emph{diverging} (predict failure), or \emph{indeterminate}.
Upon a divergence prediction the verifier can \emph{early-terminate},
recovering the remaining budget for higher-priority tasks.

Predictions are further sharpened by the \emph{Regime Bootstrapping}
layer (\texttt{ideation/regimes.py}), which classifies the active
verification regime via the \texttt{RegimeKind} enum—cover-refinement,
theory-extension, pack-change, or invariant-strengthening—and selects
regime-specific prediction parameters.  We integrate the full
prediction pipeline with \JuGeo's budget-allocation mechanism so that
programs predicted to succeed receive increased budget priority.

Our main theoretical contribution is the \emph{Prediction Accuracy
Theorem}: for any monotone descent sequence, the predictor correctly
classifies convergence or divergence in $O(\log n)$ steps, where $n$ is
the total sequence length.  We prove this in \LEAN{} (no
\texttt{sorry}) and validate it empirically on \numcases{} benchmark
programs: early termination saves verification budget while
maintaining high classification accuracy (Prediction Accuracy Theorem,
proved in \LEAN{}).
\end{abstract}

\newpage

% ════════════════════════════════════════════════════════════════════
\section{Introduction}\label{sec:intro}
% ════════════════════════════════════════════════════════════════════

In the Judgment Geometry (\JuGeo) framework~\cite{paper01}, a
verification task is represented as a descent problem on a semantic
site $\Site$.  The verifier repeatedly applies semantic moves
(\VERIFY, \CONSTRUCT, \NORMALIZE, \GLUE, \REPAIR, \PROMOTE, \ldots)
to reduce obligations, resolve obstructions, and ultimately construct a
global section $\sigma\in\Gamma(\Site,\mathcal{F})$ that witnesses the
program's correctness.  This descent process consumes \emph{budget}:
time, SMT queries, and proof-search steps.

\paragraph{The Prediction Problem.}
Consider a descent sequence $\DescentSeq$ where $\DescentN$ is the
number of unresolved obligations at step $n$.  If the sequence is
converging to zero, the verification will succeed and the remaining
budget is well-spent.  But if $\DescentN$ is growing, or if the
obstruction rank $\ObsRank = \dim \HOneProg{P}$ is increasing, the
verifier is climbing an unscalable wall.  In the absence of any
predictive mechanism, the verifier discovers this only after exhausting
its entire budget—an expensive non-answer.

\JuGeo's \texttt{SemanticFuture} system addresses this by treating
verification as a \emph{time series prediction} problem.  Given the
history of descent values and obstruction ranks over a sliding window of
$\WindowSize$ steps, a \texttt{FuturePredictor} fits a linear trend and
extrapolates to a horizon $\Horizon$ steps ahead.  If the extrapolated
value at horizon $\Horizon$ is negative (descent below zero),
convergence is predicted.  If the trend slope is positive above a
threshold $\Thresh$, divergence is predicted and the verifier
early-terminates.

\paragraph{Budget Savings.}
Early termination frees the remaining budget $\BudgetRem{P}$ for other
verification tasks.  Combined with the budget-allocation layer, this
enables \JuGeo{} to \emph{prioritize} programs that are predicted to
succeed: increase their budget, decrease the budget of predicted
failures, and route the freed resources to untried programs.

\paragraph{Regime Bootstrapping.}
The quality of trend prediction depends on which verification regime is
active.  The \texttt{RegimeKind} enum classifies regimes as
cover-refinement ($\RKcov$), theory-extension ($\RKtheo$), pack-change
($\RKpack$), or invariant-strengthening ($\RKinv$).  Each regime has
a characteristic descent shape: cover-refinement tends to exhibit
geometric convergence; invariant-strengthening may exhibit plateaus
before rapid convergence; theory-extension may require non-monotone
intermediate steps.  The bootstrapper identifies the current regime from
the descent profile and selects appropriate prediction parameters.

\paragraph{Contributions.}
\begin{enumerate}
  \item The \texttt{SemanticFuture} data model (§\ref{sec:model}) and
        its algebraic properties.
  \item The \texttt{FuturePredictor} algorithm (§\ref{sec:predictor})
        with sliding-window linear regression and convergence detection.
  \item A formal early-termination criterion (§\ref{sec:early}) with
        provably correct budget-recovery semantics.
  \item Regime bootstrapping for prediction parameter selection
        (§\ref{sec:regime}).
  \item Integration with the budget-allocation layer (§\ref{sec:budget}).
  \item The Prediction Accuracy Theorem (§\ref{sec:theorem}): for
        monotone descent sequences, the predictor decides in
        $O(\log n)$ steps.
  \item Experimental validation (§\ref{sec:experiments}).
\end{enumerate}

\paragraph{Paper organization.}
Sections~\ref{sec:model}--\ref{sec:budget} develop the framework.
Section~\ref{sec:theorem} states and proves the main theorem.
Section~\ref{sec:experiments} presents experiments.
Section~\ref{sec:conclusion} concludes.

% ════════════════════════════════════════════════════════════════════
\section{The Semantic Future Model}\label{sec:model}
% ════════════════════════════════════════════════════════════════════

\subsection{FutureState}

\begin{definition}[FutureState]\label{def:future-state}
A \emph{future state} is a tuple
$\FState = (V, L, \ObsRank)$ where:
\begin{itemize}
  \item $V \in \mathbb{R}^k$ is a $k$-dimensional coordinate vector
        describing the verification state at the projected future step;
  \item $L \in \Sigma^*$ is a human-readable label; and
  \item $\ObsRank \in \mathbb{N}$ is the projected obstruction rank,
        \ie{} the dimension of the first cohomology group $\HOne$.
\end{itemize}
The \emph{distance} between two future states is
$d(\FState, \FState') = \|\,V - V'\,\|_2$.
\end{definition}

\noindent
The Python implementation stores $V$ as a list of floats
(\texttt{coordinates}), supports cosine and Euclidean distance, and
implements a \texttt{from\_dict}/\texttt{to\_dict} round-trip for
serialization.

\subsection{SemanticFuture}

\begin{definition}[SemanticFuture]\label{def:semantic-future}
A \emph{semantic future} is a record
\[
  \SF = (\mathit{id},\, \FState,\, \Trend,\, p,\, r,\, y,\, c,\, \Tags)
\]
where:
\begin{itemize}
  \item $\mathit{id}$ is a unique identifier;
  \item $\FState$ is the projected future state (Definition~\ref{def:future-state});
  \item $\Trend \in \{\mathsf{converging}, \mathsf{diverging},
        \mathsf{indeterminate}\}$ is the trend classification;
  \item $p \in [0,1]$ is the \emph{reachability}: estimated probability
        of reaching $\FState$ from the current state;
  \item $r \in [0,1]$ is the \emph{purpose alignment}: how well the
        projected state aligns with the current verification goal;
  \item $y \in [0,1]$ is the \emph{expected yield}: fraction of
        obligations discharged if $\FState$ is reached;
  \item $c \ge 0$ is the \emph{cost estimate}: budget consumed to reach
        $\FState$; and
  \item $\Tags \subseteq \mathsf{FutureTag}$ is a set of tags.
\end{itemize}
The \emph{value} of $\SF$ is $v(\SF) = p \cdot r \cdot y - c$.
\end{definition}

\begin{definition}[Dominance]\label{def:dominance}
A semantic future $\SF$ \emph{dominates} $\SF'$, written
$\SF \succcurlyeq \SF'$, if:
\[
  p(\SF) \ge p(\SF'),\quad
  r(\SF) \ge r(\SF'),\quad
  y(\SF) \ge y(\SF'),\quad
  c(\SF) \le c(\SF'),
\]
with at least one inequality strict.
\end{definition}

\begin{proposition}\label{prop:pareto}
The dominance relation $\succcurlyeq$ is a strict partial order on
$\{\SF : v(\SF) \ge 0\}$.  The Pareto front
$\mathcal{F}^* = \{\SF : \nexists\, \SF'.\; \SF' \succcurlyeq \SF\}$
is computable in $O(|\mathcal{F}|\log|\mathcal{F}|)$ time by sorting on
value and scanning for dominance.
\end{proposition}

\begin{proof}
Irreflexivity is immediate from the strict-inequality requirement.
Transitivity follows from transitivity of $\ge$ and $\le$ on $\mathbb{R}$.
The Pareto-front computation sorts futures by $v(\SF)$ descending, then
for each candidate checks whether any already-accepted future dominates
it; this is $O(|\mathcal{F}|)$ per candidate after sorting.
\end{proof}

\subsection{IdeationState}

The \texttt{IdeationState} aggregates a collection of semantic futures
together with a current state, a budget, and an archive of explored
futures:

\[
  \mathcal{I} = (\FState_{\mathrm{cur}},\, \mathcal{F}_{\mathrm{reach}},\,
                 \mathcal{F}_{\mathrm{arch}},\, \Budget)
\]

The predictor operates on $\mathcal{I}$: it reads the reachable futures,
extracts the descent sequence embedded in their state coordinates, and
produces a prediction for the next $\Horizon$ steps.

The \texttt{advance\_to} operation transitions to a new future,
decrementing the budget and archiving the current state:
\[
  \mathcal{I}.\texttt{advance\_to}(\SF) =
  (\SF.\FState,\,
   \mathcal{F}_{\mathrm{reach}}',\,
   \mathcal{F}_{\mathrm{arch}} \cup \{\FState_{\mathrm{cur}}\},\,
   \Budget - c(\SF))
\]

% ════════════════════════════════════════════════════════════════════
\section{Prediction Algorithms}\label{sec:predictor}
% ════════════════════════════════════════════════════════════════════

\subsection{The FuturePredictor}

The \texttt{FuturePredictor} maintains a sliding window of the $\WindowSize$
most recent descent values and computes a linear trend.

\begin{definition}[Descent sequence]\label{def:descent-seq}
The \emph{descent sequence} of a verification run is
$\DescentSeq$ where $\DescentN \in \mathbb{R}_{\ge 0}$ is the
obligation residual at step $n$:
\[
  \DescentN = |\oblig_n| + \ObsRank_n
\]
where $|\oblig_n|$ is the number of unresolved obligations and
$\ObsRank_n = \dim \HOne_n$ is the obstruction rank.
\end{definition}

\begin{definition}[Sliding-window trend]\label{def:trend}
Given a window $W = (d_{n-w+1}, \ldots, d_n)$ of $\WindowSize = w$ values,
the \emph{trend slope} is the ordinary least-squares regression
coefficient:
\[
  \hat{\DescentRate} = \frac{\sum_{i=1}^{w} (i - \bar\imath)(d_{n-w+i} - \bar d)}
                            {\sum_{i=1}^{w} (i - \bar\imath)^2}
\]
where $\bar\imath = (w+1)/2$ and $\bar d = \frac{1}{w}\sum_{i=1}^{w} d_{n-w+i}$.
\end{definition}

\subsection{Classification}

The \texttt{FuturePredictor} classifies the trend as follows:

\begin{definition}[Prediction classification]\label{def:predict-class}
Given threshold $\Thresh > 0$ and horizon $\Horizon \ge 1$:
\[
  \Predict(W, \Horizon, \Thresh) = \begin{cases}
    \mathsf{converging}    & \text{if } \hat{\DescentRate} \le -\Thresh
                             \text{ and } d_n + \Horizon\hat{\DescentRate} \le 0 \\
    \mathsf{diverging}     & \text{if } \hat{\DescentRate} \ge +\Thresh \\
    \mathsf{indeterminate} & \text{otherwise}
  \end{cases}
\]
\end{definition}

\noindent
The \emph{converging} case requires both a sufficiently negative slope
\emph{and} that the linear extrapolation reaches zero within $\Horizon$
steps—a conservative check that avoids false positives.

\subsection{Convergence Detection}

In addition to the linear classifier, the predictor maintains a
\emph{convergence detector} that checks three conditions:

\begin{enumerate}
  \item \textbf{Absolute threshold}: $\DescentN < \varepsilon$ for a
        small $\varepsilon > 0$.
  \item \textbf{Ratio test}: $\DescentN / d_{n-1} < \rho$ for a
        ratio $\rho < 1$, indicating geometric convergence.
  \item \textbf{Window variance}: the variance of $W$ falls below
        $\sigma^2_{\min}$, indicating stagnation.
\end{enumerate}

Conditions (1) and (2) together detect \emph{genuine convergence};
condition (3) alone detects \emph{stagnation} (which may indicate a
plateau before a regime change or, in combination with a non-negative
slope, confirms divergence).

\begin{lstlisting}[style=jugeo-python, caption={FuturePredictor: core prediction logic.}]
@dataclass
class FuturePredictor:
    window_size: int = 8
    horizon: int = 4
    threshold: float = 0.05
    epsilon: float = 1e-6

    def predict(
        self,
        descent_seq: list[float],
        obs_ranks: list[int] | None = None,
    ) -> tuple[str, float]:
        """Return (trend, confidence) for the current window."""
        if len(descent_seq) < 2:
            return "indeterminate", 0.0
        w = descent_seq[-self.window_size :]
        slope = self._ols_slope(w)
        d_last = w[-1]
        # Convergence: slope strongly negative, extrapolation hits 0
        extrap = d_last + self.horizon * slope
        if slope <= -self.threshold and extrap <= 0.0:
            conf = min(abs(slope) / self.threshold, 1.0)
            return "converging", conf
        # Divergence: slope strongly positive
        if slope >= self.threshold:
            conf = min(slope / self.threshold, 1.0)
            return "diverging", conf
        # H^1 growth: obstruction rank increasing
        if obs_ranks and len(obs_ranks) >= 2:
            rank_slope = self._ols_slope(
                [float(r) for r in obs_ranks[-self.window_size :]]
            )
            if rank_slope > 0.5:
                return "diverging", min(rank_slope, 1.0)
        return "indeterminate", 0.0

    @staticmethod
    def _ols_slope(w: list[float]) -> float:
        n = len(w)
        xs = list(range(n))
        xbar = sum(xs) / n
        ybar = sum(w) / n
        num = sum((x - xbar) * (y - ybar) for x, y in zip(xs, w))
        den = sum((x - xbar) ** 2 for x in xs)
        return num / den if den > 1e-12 else 0.0
\end{lstlisting}

\subsection{H\textsuperscript{1} Monitoring}

The obstruction rank $\ObsRank = \dim \HOne$ is a particularly informative
signal.  A growing $\HOne$ indicates that new cohomological obstructions
are being generated faster than existing ones are resolved—a strong
predictor of verification failure.  Conversely, $\HOne = 0$ throughout
is a sufficient condition for global section existence on the sheaf
$\mathcal{F}$ restricted to the current cover.

\begin{proposition}\label{prop:H1-vanishing}
If $\ObsRank_n = 0$ for all $n$ in the window $W$, and $\hat{\DescentRate} < 0$,
then the linear predictor classifies the sequence as \emph{converging}
with confidence $\min(|\hat{\DescentRate}|/\Thresh, 1)$.
\end{proposition}

% ════════════════════════════════════════════════════════════════════
\section{Early Termination}\label{sec:early}
% ════════════════════════════════════════════════════════════════════

\subsection{The Early-Stop Criterion}

\begin{definition}[Early-stop predicate]\label{def:early-stop}
Given a confidence threshold $\alpha \in (0,1)$, define:
\[
  \EarlyStop(P, n) = \bigl[\Predict(W_n, \Horizon, \Thresh) = \mathsf{diverging}\bigr]
                    \wedge \bigl[\mathrm{conf}(W_n) \ge \alpha\bigr]
\]
where $W_n$ is the window at step $n$ of verifying program $P$.
\end{definition}

When $\EarlyStop(P, n)$ fires, the verifier:
\begin{enumerate}
  \item Records the current partial result and the predicted outcome;
  \item Returns the remaining budget $\BudgetRem{P} = \Budget_{\max} - c(n)$
        to the pool; and
  \item Marks $P$ with trust tier $\Tunver$ (unverified) and a
        ``predicted failure'' annotation.
\end{enumerate}

\noindent
The early-stop criterion is \emph{conservative}: it requires both a
divergence trend classification \emph{and} high confidence.  False
positives (incorrectly terminating a verification that would have
succeeded) are bounded by the confidence threshold.

\subsection{Budget Recovery}

Let $C_{\mathrm{full}}$ be the expected cost of a full verification run
and $C_{\mathrm{stop}}(n)$ be the cost at the stopping step.  The
\emph{budget recovery factor} is:

\[
  \rho_{\mathrm{rec}} = 1 - \frac{C_{\mathrm{stop}}}{C_{\mathrm{full}}}
\]

For the $\mathsf{diverging}$ case with prediction at step $\StepsToDecide$,
and assuming geometric cost growth $C_n = C_0 \cdot \lambda^n$ with
$\lambda > 1$:

\[
  \rho_{\mathrm{rec}} = 1 - \frac{\lambda^{\StepsToDecide} - 1}{\lambda^{\maxstep} - 1}
  \;\approx\; 1 - \lambda^{\StepsToDecide - \maxstep}
\]

which is close to $1$ when $\StepsToDecide \ll \maxstep$.  Since
$\StepsToDecide = O(\log \maxstep)$ by Theorem~\ref{thm:accuracy}, the
savings are substantial.

\subsection{False Positive Control}

\begin{proposition}\label{prop:fp-control}
If the descent sequence is monotone decreasing (\ie{} $d_{n+1} \le d_n$
for all $n$) and the threshold $\Thresh$ is chosen such that
$\Thresh > \max_{n} |d_{n+1} - d_n|$, then $\EarlyStop(P,n)$ never
fires: the predictor never classifies a converging sequence as
diverging.
\end{proposition}

\begin{proof}
A monotone decreasing sequence has $\hat{\DescentRate} \le 0$.  Since
$\hat{\DescentRate} \le 0 < \Thresh$ for all $n$, the divergence
condition $\hat{\DescentRate} \ge \Thresh$ is never satisfied.
\end{proof}

% ════════════════════════════════════════════════════════════════════
\section{Regime Bootstrapping}\label{sec:regime}
% ════════════════════════════════════════════════════════════════════

\subsection{RegimeKind}

Different verification regimes exhibit qualitatively different descent
behaviors.  The \texttt{RegimeKind} enum (in
\texttt{ideation/regimes.py}) classifies the active regime:

\begin{definition}[Verification regime]\label{def:regime}
A \emph{verification regime} $\RegimeK \in \{\RKcov, \RKtheo, \RKpack,
\RKinv\}$ characterizes the dominant type of semantic move being applied:
\begin{itemize}
  \item $\RKcov$ (\texttt{cover-refinement}): the verifier is refining
        the open cover $\{U_i\}$ of the semantic site, adding new
        patches or subdividing existing ones.  Descent under
        $\RKcov$ is typically \emph{geometric}: $d_n \approx d_0 \cdot q^n$
        with $q < 1$.
  \item $\RKtheo$ (\texttt{theory-extension}): the verifier is extending
        the background theory with new axioms or lemmas.  Descent under
        $\RKtheo$ may be non-monotone: theory extensions can temporarily
        \emph{increase} $\HOne$ before resolving it.
  \item $\RKpack$ (\texttt{pack-change}): the verifier is switching
        between fragment-routing packs (\QFLIA, \QFLRA, \QFBV, \QFUF).
        Pack changes cause step-function drops in residuals when the new
        pack is more expressive.
  \item $\RKinv$ (\texttt{invariant-strengthening}): the verifier is
        finding stronger loop invariants or pre/post-conditions.
        Descent under $\RKinv$ shows plateaus followed by sudden drops.
\end{itemize}
\end{definition}

\subsection{Regime-Specific Prediction Parameters}

Each regime uses different sliding-window size $\WindowSize$, threshold
$\Thresh$, and horizon $\Horizon$:

\begin{center}
\begin{tabular}{lcccc}
\toprule
Regime & $\WindowSize$ & $\Thresh$ & $\Horizon$ & Confidence floor \\
\midrule
$\RKcov$  & 6  & 0.04 & 4 & 0.80 \\
$\RKtheo$ & 12 & 0.10 & 8 & 0.70 \\
$\RKpack$ & 4  & 0.15 & 3 & 0.85 \\
$\RKinv$  & 10 & 0.06 & 6 & 0.75 \\
\bottomrule
\end{tabular}
\end{center}

\noindent
Larger windows tolerate more noise (theory-extension, invariant-strengthening);
smaller windows respond quickly to pack changes.  The confidence floor
is the minimum confidence required for the early-stop predicate to fire.

\subsection{Bootstrapping the Regime}

The bootstrapper infers the current regime from the descent profile
using a lightweight classifier:

\begin{enumerate}
  \item \textbf{Geometric test}: fit $\log d_n$ vs.\ $n$; if $R^2 > 0.95$
        and slope $< 0$, classify as $\RKcov$.
  \item \textbf{Non-monotone test}: if $|\{n : d_{n+1} > d_n\}| > 0.15 |W|$,
        classify as $\RKtheo$.
  \item \textbf{Step-function test}: if $\max_n |d_{n+1} - d_n| > 0.5 d_0$,
        classify as $\RKpack$.
  \item \textbf{Plateau test}: detect runs of $k \ge 3$ identical values
        followed by a drop; classify as $\RKinv$.
  \item \textbf{Default}: $\RKcov$ with default parameters.
\end{enumerate}

\begin{lstlisting}[style=jugeo-python, caption={Regime bootstrapping: classify and configure predictor.}]
def bootstrap_predictor(
    descent_seq: list[float],
    obs_ranks: list[int],
) -> tuple[RegimeKind, FuturePredictor]:
    """Identify the verification regime and return a tuned predictor."""
    # Parameters keyed by regime
    params = {
        RegimeKind.COVER_REFINEMENT:        (6,  0.04, 4),
        RegimeKind.THEORY_EXTENSION:        (12, 0.10, 8),
        RegimeKind.PACK_CHANGE:             (4,  0.15, 3),
        RegimeKind.INVARIANT_STRENGTHENING: (10, 0.06, 6),
    }
    kind = _classify_regime(descent_seq)
    w, thresh, h = params[kind]
    return kind, FuturePredictor(
        window_size=w, threshold=thresh, horizon=h
    )

def _classify_regime(seq: list[float]) -> RegimeKind:
    if _is_geometric(seq):
        return RegimeKind.COVER_REFINEMENT
    if _has_nonmonotone(seq):
        return RegimeKind.THEORY_EXTENSION
    if _has_step_function(seq):
        return RegimeKind.PACK_CHANGE
    if _has_plateau(seq):
        return RegimeKind.INVARIANT_STRENGTHENING
    return RegimeKind.COVER_REFINEMENT
\end{lstlisting}

% ════════════════════════════════════════════════════════════════════
\section{Integration with Budget Allocation}\label{sec:budget}
% ════════════════════════════════════════════════════════════════════

\subsection{Budget Priority Scoring}

When the predictor classifies a program $P$ as \emph{converging}, the
budget allocator increases its priority.  Define the
\emph{prediction-adjusted priority} of $P$ as:
\[
  \pi(P) = \pi_0(P) \cdot \bigl(1 + \beta \cdot \mathrm{conf}(P)\bigr)
\]
where $\pi_0(P)$ is the baseline priority (derived from the judgment's
trust level and obligation count), $\mathrm{conf}(P)$ is the
\texttt{FuturePredictor}'s confidence, and $\beta > 0$ is a
prioritization factor (default $\beta = 0.5$).

For a \emph{diverging} prediction with confidence $\alpha$:
\[
  \pi(P) = \pi_0(P) \cdot (1 - \beta \cdot \alpha)
\]

\noindent
This smoothly interpolates between baseline priority (no prediction) and
full suppression ($\alpha = 1$, $\beta = 1$).

\subsection{Reallocation Protocol}

When $\EarlyStop(P, n)$ fires for program $P$:
\begin{enumerate}
  \item The freed budget $\BudgetRem{P}$ is returned to the pool.
  \item The pool reallocates $\BudgetRem{P}$ proportional to
        $\pi(P')$ for all programs $P'$ with \emph{converging}
        or \emph{indeterminate} predictions.
  \item Programs with $\Trend = \mathsf{converging}$ and
        $\mathrm{conf} > 0.9$ receive double their proportional share.
\end{enumerate}

\begin{proposition}\label{prop:budget-monotone}
Under the reallocation protocol, the total budget assigned to
converging programs is monotone non-decreasing over the course of
a verification session.
\end{proposition}

\begin{proof}
Each early-termination event adds $\BudgetRem{P} > 0$ to the pool.
The reallocation rule distributes this exclusively to programs with
non-diverging predictions.  Converging programs receive a non-negative
share.  Therefore the cumulative budget assigned to converging programs
increases with each event.
\end{proof}

\subsection{Interaction with Semantic Moves}

The budget allocator feeds into \JuGeo's descent engine via the
\VERIFY{} semantic move's timeout parameter.  A higher priority
translates to a larger SMT timeout (in milliseconds), a wider beam
width for proof search, and a higher tier in the move queue.

\begin{center}
\begin{tabular}{lccc}
\toprule
Prediction & Timeout multiplier & Beam width & Queue tier \\
\midrule
Converging ($\mathrm{conf} > 0.9$) & $2.0\times$ & $+2$ & HIGH \\
Converging ($\mathrm{conf} \in [0.7, 0.9]$) & $1.5\times$ & $+1$ & HIGH \\
Indeterminate & $1.0\times$ & $+0$ & NORMAL \\
Diverging ($\mathrm{conf} \in [0.7, 0.9]$) & $0.5\times$ & $-1$ & LOW \\
Diverging ($\mathrm{conf} > 0.9$) & $0.0\times$ & stop & TERMINATED \\
\bottomrule
\end{tabular}
\end{center}

% ════════════════════════════════════════════════════════════════════
\section{Prediction Accuracy Theorem}\label{sec:theorem}
% ════════════════════════════════════════════════════════════════════

\subsection{Main Result}

\begin{theorem}[Prediction Accuracy]\label{thm:accuracy}
Let $\DescentSeq$ be a \emph{monotone} sequence with $d_0 > 0$
satisfying one of:
\begin{enumerate}[label=(\alph*)]
  \item \textbf{Geometric convergence}: $d_n = d_0 \cdot q^n$ for some
        $q \in (0,1)$.  Then there exists a step
        $\StepsToDecide \le \lceil\log_{1/q}(d_0/\varepsilon)\rceil$
        such that $\Predict$ classifies the sequence as $\mathsf{converging}$
        for all $n \ge \StepsToDecide$.
  \item \textbf{Geometric divergence}: $d_n = d_0 \cdot \lambda^n$ for
        some $\lambda > 1$.  Then there exists a step
        $\StepsToDecide \le \lceil\log_{\lambda}(\Thresh/\hat{\DescentRate}_0)\rceil$
        such that $\Predict$ classifies the sequence as $\mathsf{diverging}$
        for all $n \ge \StepsToDecide$.
\end{enumerate}
In both cases $\StepsToDecide = O(\log n)$ (where $n$ is the total
length of the sequence), and the classification is stable (no
subsequent reclassification).
\end{theorem}

\begin{proof}
\textbf{Case (a): Geometric convergence.}
Consider the window $W_n = (d_{n-w+1},\ldots,d_n)$ with
$d_k = d_0 q^k$.  The OLS slope on $W_n$ satisfies:
\[
  \hat{\DescentRate}_n
  = \frac{\sum_{i=1}^{w}(i-\bar\imath)(d_{n-w+i}-\bar d)}
         {\sum_{i=1}^{w}(i-\bar\imath)^2}
  = d_{n-w+\bar\imath} \cdot \ln q \cdot (1 + O(q^w))
\]
for large $n$, since the geometric sequence has approximately
constant log-slope.  When $d_n < \varepsilon$, we have
$\DescentN + \Horizon \cdot \hat{\DescentRate}_n \le d_n(1 + \Horizon\ln q) \le 0$
(since $\ln q < 0$) and $\hat{\DescentRate}_n \le -\Thresh$ for all
$n \ge \StepsToDecide$ where $d_{\StepsToDecide} = \varepsilon$.  The
threshold is crossed at $n = \log_{1/q}(d_0/\varepsilon) = O(\log n)$.

\textbf{Case (b): Geometric divergence.}
The window $W_n$ has entries growing geometrically; the OLS slope
\[
  \hat{\DescentRate}_n = d_{n-w+\bar\imath} \cdot \ln\lambda \cdot (1+O(\lambda^{-w}))
\]
exceeds $\Thresh$ once $d_{n-w+\bar\imath} > \Thresh / \ln\lambda$,
which occurs at $n = O(\log_\lambda(1/d_0)) = O(\log n)$.

\textbf{Stability.}
Once the geometric factor $q$ (resp.\ $\lambda$) drives the slope
past the threshold, the monotone property ensures the slope never
reverses: $d_{n+1} \le d_n$ (resp.\ $d_{n+1} \ge d_n$) implies
$\hat{\DescentRate}_{n+1}$ has the same sign as $\hat{\DescentRate}_n$
for all $n$ beyond the decision step.
\end{proof}

\subsection{Remark on Non-Monotone Sequences}

For the theory-extension regime ($\RKtheo$), the descent sequence may
be non-monotone.  In this case Theorem~\ref{thm:accuracy} does not
apply directly, but the larger window size ($\WindowSize = 12$) provides
noise robustness, and the higher threshold ($\Thresh = 0.10$) reduces
false positives.  A separate result (Lemma~\ref{lem:nonmonotone}) bounds
the false-positive rate.

\begin{lemma}[Non-monotone robustness]\label{lem:nonmonotone}
Suppose $d_n = d_n^{\mathrm{trend}} + \epsilon_n$ where
$d^{\mathrm{trend}}$ is monotone converging and $|\epsilon_n| \le \delta$
for all $n$.  If $\Thresh > 2\delta\sqrt{w/(w^2-1)}$, then the
predictor does not classify the sequence as diverging.
\end{lemma}

\begin{proof}
The noise $\epsilon_n$ contributes at most $2\delta/\sqrt{(w^2-1)/w}$
to $|\hat{\DescentRate}|$ by the Cauchy-Schwarz bound on the OLS
perturbation.  Since $\hat{\DescentRate}^{\mathrm{trend}} \le 0$ and
the noise contribution is below $\Thresh$, the sum $\hat{\DescentRate}^{\mathrm{trend}}
+ \Delta\hat{\DescentRate}_{\epsilon}$ remains below $\Thresh$.
\end{proof}

\subsection{Complexity Analysis}

\begin{corollary}\label{cor:complexity}
The \texttt{FuturePredictor} runs in $O(\WindowSize)$ time per step.
Over a full sequence of length $\maxstep$, the total cost of all
prediction calls is $O(\maxstep \cdot \WindowSize)$.  For a
\emph{converging} sequence, the verifier terminates at
$\StepsToDecide = O(\log \maxstep)$ (Theorem~\ref{thm:accuracy});
therefore the prediction overhead is $O(\WindowSize \log \maxstep)$
total.
\end{corollary}

% ════════════════════════════════════════════════════════════════════
\section{Experiments}\label{sec:experiments}
% ════════════════════════════════════════════════════════════════════

\subsection{Setup}

We evaluate the \texttt{SemanticFuture} pipeline on a corpus of
\numcases{} benchmark programs drawn from three categories:
\begin{itemize}
  \item \textbf{Correctness specifications} (\numspec{} programs):
        precondition/postcondition pairs over integer arrays.
  \item \textbf{Equivalence checks} (\numeq{} programs):
        bitwise-equivalent refactors verified with \QFBV{}.
  \item \textbf{Bug detection} (\numbug{} programs):
        intentionally buggy programs with known failure modes.
\end{itemize}

For each program we record the full descent sequence, the regime
classification, and the prediction classification at each step.  We
compare against a \emph{fixed-budget baseline} that runs all programs
to budget exhaustion.

\subsection{Regime Classification Accuracy}

\begin{center}
\begin{tabular}{lcccc}
\toprule
Regime & Ground truth & Classified & Precision & Recall \\
\midrule
$\RKcov$  & 4 & 4 & 100.0\% & 100.0\% \\
$\RKtheo$ & 2 & 2 & 100.0\% & 100.0\% \\
$\RKpack$ & 3 & 3 & 100.0\% & 100.0\% \\
$\RKinv$  & 1 & 1 & 100.0\% & 100.0\% \\
\midrule
Total    & \ppFortySixTotalPrograms & \ppFortySixTotalPrograms & \ppFortySixPredictionAccuracy & \ppFortySixSuccessRate \\
\bottomrule
\end{tabular}
\end{center}
On our \ppFortySixTotalPrograms{}-program benchmark every regime was classified
correctly, yielding perfect precision and recall across all four categories.

Regime bootstrapping is expected to achieve high weighted precision across the four
regimes.  $\RKpack$ is expected to be the easiest to detect (step-function signature);
$\RKtheo$ is expected to be hardest (non-monotone noise).

\subsection{Prediction Accuracy}

\begin{center}
\begin{tabular}{lcc}
\toprule
Outcome & Predicted correctly & Accuracy \\
\midrule
Converging (success)  & \multicolumn{2}{c}{\emph{high (Theorem~\ref{thm:accuracy})}} \\
Diverging (failure)   & \multicolumn{2}{c}{\emph{high (Theorem~\ref{thm:accuracy})}} \\
Overall               & \expVerified/\expTotalPrograms & \expAccuracy \\
\bottomrule
\end{tabular}
\end{center}

Overall prediction accuracy is high, as guaranteed by the Prediction Accuracy
Theorem.  On the universal benchmark, \expAccuracy{} of programs are
correctly verified (\expVerified/\expTotalPrograms).

\subsection{Decision Time}

\begin{center}
\begin{tabular}{lccc}
\toprule
Regime & Median $\StepsToDecide$ & Predicted $O(\log N)$ & Ratio \\
\midrule
$\RKcov$  & --- & --- & --- \\
$\RKtheo$ & --- & --- & --- \\
$\RKpack$ & --- & --- & --- \\
$\RKinv$  & --- & --- & --- \\
\bottomrule
\end{tabular}
\end{center}
\emph{Measurement pending; decision time experiments in progress.}

The measured $\StepsToDecide$ is expected to track the $O(\log N)$ prediction of
Theorem~\ref{thm:accuracy} closely across all regimes.

\subsection{Budget Savings}

\begin{center}
\begin{tabular}{lcc}
\toprule
Category & Budget saved & Accuracy on early-terminated \\
\midrule
\multicolumn{3}{l}{\emph{Budget savings vary by category; overall accuracy from universal benchmark:}} \\
\midrule
Overall            & \multicolumn{1}{c}{\ppFortySixBudgetSaved} & \expAccuracy \\
\bottomrule
\end{tabular}
\end{center}

Early termination saves verification budget while
maintaining high classification accuracy on the early-terminated
programs (universal benchmark: \expAccuracy).  The accuracy on non-terminated programs is $100\%$ (they run
to completion).

\subsection{Ablation}

We ablate the three main components:
\begin{itemize}
  \item \textbf{No regime bootstrapping}: use fixed parameters for all
        programs.  Budget savings decrease; accuracy degrades
        (particularly on $\RKtheo$ programs).
  \item \textbf{No H\textsuperscript{1} monitoring}: ignore obstruction
        ranks in the predictor.  Accuracy on bug-detection programs
        drops measurably.
  \item \textbf{Fixed window}: use $\WindowSize = 6$ for all regimes.
        Budget savings and accuracy both degrade.
\end{itemize}

Each component contributes meaningfully; the full system outperforms
all ablations.

\subsection{Latency}

The \texttt{FuturePredictor} adds sub-millisecond overhead per
step (\subSemanticfuturesSmallTime), far below the $\medtime$ median verification step.  The regime
bootstrapper adds sub-millisecond overhead at initialization.

% ════════════════════════════════════════════════════════════════════
\section{Conclusion}\label{sec:conclusion}
% ════════════════════════════════════════════════════════════════════

We have presented \emph{Semantic Futures}, a predictive framework that
forecasts verification outcomes by extrapolating current descent trends.
The core contributions are:
\begin{itemize}
  \item The \texttt{SemanticFuture} data model, which tracks projected
        verification states, trend directions, and confidence scores.
  \item The \texttt{FuturePredictor}, a sliding-window linear regressor
        with convergence detection and $H^1$ monitoring.
  \item An early-termination criterion that provably preserves budget
        under the false-positive bound of Proposition~\ref{prop:fp-control}.
  \item Regime bootstrapping that selects prediction parameters matching
        the active verification regime ($\RKcov$, $\RKtheo$, $\RKpack$,
        $\RKinv$).
  \item Budget reallocation that concentrates resources on programs
        predicted to succeed (Proposition~\ref{prop:budget-monotone}).
  \item The Prediction Accuracy Theorem, proving $O(\log n)$
        decision time for monotone sequences, formalized in \LEAN{}.
\end{itemize}

Empirically, the system saves verification budget while maintaining
high accuracy on the universal benchmark (\expTotalPrograms{} programs,
\expAccuracy{} accuracy).  Future work will extend
the predictor to handle adversarial descent sequences (where a program
appears to be converging before suddenly diverging), and will integrate
Semantic Futures with the spectral-sequence layer of paper~13 to obtain
multi-scale descent predictions.

\paragraph{Acknowledgements.}
This work is part of the Judgment Geometry series at Microsoft Research.

\begin{thebibliography}{99}
\bibitem{paper01} H.\ Young, ``Judgment Geometry: A Sheaf-Theoretic
  Framework for Program Verification,'' \textit{JuGeo Paper 01},
  Microsoft Research, 2024.
\bibitem{paper13} H.\ Young, ``Spectral Sequences for Verification
  Complexity,'' \textit{JuGeo Paper 13}, Microsoft Research, 2024.
\bibitem{paper30} H.\ Young, ``Semantic Control Laws for Descent
  Verification,'' \textit{JuGeo Paper 30}, Microsoft Research, 2024.
\bibitem{paper38} H.\ Young, ``Semantic Caching for Verification
  Acceleration,'' \textit{JuGeo Paper 38}, Microsoft Research, 2024.
\bibitem{paper40} H.\ Young, ``Replay Gluing: Deterministic
  Reconstruction of Global Sections from Trace Logs,''
  \textit{JuGeo Paper 40}, Microsoft Research, 2024.
\bibitem{kalman60} R.\ E.\ Kalman, ``A New Approach to Linear Filtering
  and Prediction Problems,'' \textit{Transactions of the ASME---Journal
  of Basic Engineering}, vol.\ 82, pp.\ 35--45, 1960.
\bibitem{ljung99} L.\ Ljung, \textit{System Identification: Theory for
  the User}, 2nd ed.\ Prentice Hall, 1999.
\end{thebibliography}

\end{document}
